\section{多中心的秩序与北方边界的价格}
\label{sec:regional-order-northern-price}

907年后梁定都开封，前蜀在成都称帝，岐与吴继续使用唐天祐年号，吴越、闽和岭南则以王号、节度使名义或新帝号组织地方。不同称号并非对“统一”简单的赞成或反对：它们为军府、地方官僚、港口与腹地之间提供可用的外交和任命语言。918年高丽建国也提醒人们，东亚的政权重组不以开封为唯一尺度；中原史书的五代年序是重要骨架，却不能替代区域政治的时间。%
\footnote{后梁、前蜀、岐、吴、吴越、闽、南汉与高丽的建立或定位见新增账本 LT-907-081 至 LT-918-091，依据《旧五代史》《新五代史》相关本纪、世家及《高丽史》早期记载。王号与年号只能确证制度语言，不能量化各政权对山地、港口或乡村的控制。}

923年后唐灭后梁、925年灭前蜀，926年庄宗即在兵变中被杀；934年后蜀重建，937年南唐取吴而立，945年又灭闽。战役和受禅改变了都城与军队的归属，却没有把巴蜀、江淮或福建转化为同样密度的行政空间。前蜀覆亡后仍能有后蜀，正说明接收留下的兵粮、任官与地方关系可以成为新的建国资源；南唐的扩张同样不能使闽岭各地立即共享一套税役和军政。%
\footnote{后唐、前后蜀、吴与南唐、闽的转折见 LT-923-092、LT-925-094、LT-926-095、LT-934-098、LT-937-100、LT-945-104，依据《旧五代史》《新五代史》本纪、前后蜀和南唐世家与《资治通鉴》。征服者的军功叙事不能代替地方接管和人口流动的连续证据。}

936年石敬瑭借契丹援助建后晋，并承诺燕云十六州；942年石重贵改变称臣关系，至946年契丹入开封，947年又北撤，后汉在太原与中原间建立。所谓“燕云之失”并非一条抽象国界的移动，而是道路、州城、牧地、互市与长期军费被置入一场建国交换。割让文书原件不存，后世“儿皇帝”的语言也有强烈道德色彩；但北方边防的供给与指挥从此不能按后唐时期的条件处理，是更可靠的结论。%
\footnote{后晋、契丹南下、后晋覆亡、辽军北撤和后汉建立见 LT-936-099、LT-941-101 至 LT-947-107，依据《旧五代史·晋纪》《辽史·太宗纪》《新五代史》与《资治通鉴》。燕云范围、使节礼仪和撤军原因应保留材料差异。}

951年郭威建后周、北汉同时在太原立国并求援于辽；954年高平之战、955年整顿寺院、956--958年南征与接收南唐江北，使后周得以把北防、铜料与漕运重新联结。959年北伐未及幽云而止，柴荣去世；960年陈桥受禅则把开封中枢和禁军交给宋。这里的连续性并非预告了统一必然到来，而是显示宋继承的是一套仍要面对辽、北汉和南方诸国的军政问题。%
\footnote{后周、北汉、高平、寺院整顿、南唐战争及陈桥受禅见 LT-951-109 至 LT-960-118，依据《旧五代史·周纪》《五代会要》《辽史》《宋史·太祖本纪》及《资治通鉴》。寺院废并、割地州数、北伐范围和陈桥预谋都不宜由单一正统叙事断定。}
